% Appendix: the practice exam, from exam/practice_exam.md. The questions only; the solutions are
% the next appendix, as they are a separate file, so that the paper can be sat without them in
% view. The source's title names the class it was first written for; the book leaves the name out.

\chapter{Övningsprov}
\label{app:practice}

Övningsprovet har samma form och samma poäng som det riktiga provet, som bilaga~\ref{app:exam}
beskriver. Gör det under två timmar, som på provet, och rätta det sedan mot lösningsförslagen i
bilaga~\ref{app:practice:answers}.

\begin{itemize}
  \item \textbf{Tid:} 2 timmar.
  \item \textbf{Hjälpmedel:} penna, suddgummi och linjal, och det utdelade referensbladet med
    AVR-instruktionerna, bilaga~\ref{app:avr}. Ingen miniräknare.
  \item \textbf{Poäng:} 40. \textbf{G} kräver 20 poäng och \textbf{VG} 30 poäng.
\end{itemize}

\noindent\textbf{Instruktioner:}
\begin{itemize}
  \item Börja varje fråga på ett nytt papper, och skriv ditt namn på varje papper.
  \item Visa hur du räknar. Metoden ger poäng även när svaret blir fel.
  \item Skriv alltid ut talsystemet: \code{0x2C}, \code{0010 1100} eller 44.
  \item Assemblerfrågorna gäller ATmega328P med klockfrekvensen 16~MHz.
\end{itemize}

% ------------------------------------------------------------------------------------------
\question{Talsystem och koder}{8}
\label{practice:q:1}

\qpart{a} Omvandla det binära talet \code{1101 0110} till decimalt och hexadecimalt. \worth{2}

\qpart{b} Omvandla det decimala talet 173 till binärt och hexadecimalt. \worth{2}

\qpart{c} Talen \code{0xB4} och \code{0x6A} adderas i ett 8-bitars register. Ställ upp additionen
binärt, och ange resultatet i registret hexadecimalt. Blir det en minnessiffra ut ur registret?
\worth{2}

\qpart{d} Skriv talet 59 i BCD. Ange också de två ASCII-koder som texten \code{59} består av.
\worth{2}

% ------------------------------------------------------------------------------------------
\question{Grindar och boolesk algebra}{8}
\label{practice:q:2}

\qpart{a} Gör sanningstabellen för \code{X = A'B + BC'}. \worth{2}

\needspace{6\baselineskip}
\qpart{b} Förenkla uttrycket algebraiskt så långt det går, och ange vilken lag du använder i varje
steg. Vilken enda grind beräknar det förenklade uttrycket? \worth{3}

\begin{console}
X = AB + AB' + A'B
\end{console}

\qpart{c}\leavevmode
\begin{enumerate}
  \item Skriv \code{(A + B)'} utan parentes. \worth{1}
  \item Visa hur funktionen \code{A + B} kan byggas av enbart NAND-grindar med två ingångar. Rita
    nätet, och visa med boolesk algebra att det stämmer. \worth{2}
\end{enumerate}

% ------------------------------------------------------------------------------------------
\question{Karnaughdiagram}{6}
\label{practice:q:3}

En funktion \code{X} av de fyra variablerna A, B, C och D är 1 för mintermerna 0, 2, 5, 7, 8, 10,
13 och 15, och 0 för alla andra.

\qpart{a} Rita Karnaughdiagrammet, med \code{AB} på raderna och \code{CD} på kolumnerna, och fyll i
ettorna. \worth{2}

\qpart{b} Gruppera ettorna minimalt, och skriv det minimerade uttrycket för \code{X} på SP-form.
\worth{3}

\qpart{c} Uttrycket beror bara på två av variablerna. Vilka två, och vilken grind med två ingångar
är \code{X}? \worth{1}

% ------------------------------------------------------------------------------------------
\question{Analys av ett kontaktnät}{6}
\label{practice:q:4}

Kontaktnätet i figur~\ref{practice:fig:contacts} styr lampan H1 med tre tryckknappar. Tryckknappen
S1 har två kontaktblock, som båda manövreras av samma knapp. S1 och S2 är slutande, S3 är brytande.

\bookfigure{exam_practice_contacts}{Kontaktnätet som styr lampan H1.}{practice:fig:contacts}

\qpart{a} Skriv det booleska uttrycket för H1. \worth{2}

\qpart{b} Gör sanningstabellen för H1, med S1, S2 och S3 som ingångar. \worth{2}

\qpart{c} Förenkla uttrycket, och rita ett kontaktnät som gör samma sak med färre kontakter.
\worth{2}

% ------------------------------------------------------------------------------------------
\question{Mikrodatorns uppbyggnad}{4}
\label{practice:q:5}

\qpart{a} Rita en mikrodators blockschema, med CPU (ALU, register, styrenhet och programräknare),
programminne, dataminne, in- och utportar och bussar. Skriv en mening om vad varje block gör.
\worth{2}

\qpart{b} Vad innehåller programräknaren, och vad händer med den när ett \code{rcall} utförs?
\worth{1}

\qpart{c} I vilket minne i ATmega328P ligger programmet, och i vilket ligger stacken? Vilket av dem
behåller sitt innehåll när strömmen bryts? \worth{1}

% ------------------------------------------------------------------------------------------
\question{Assemblerprogrammering}{8}
\label{practice:q:6}

\qpart{a} Programmet nedan körs från början. Ange värdena i \code{r16}, \code{r17} och \code{r18}
hexadecimalt när det har kommit till \code{end}, och värdet på flaggan \code{C} efter \code{add}.
Förklara varför \code{r18} inte blev 300. \worth{3}

\begin{asmcode}
    ldi r16, 0xB4
    ldi r17, 0x0F
    and r17, r16
    lsl r16
    ldi r18, 200
    ldi r19, 100
    add r18, r19
end:
    rjmp end
\end{asmcode}

\qpart{b} Programmet nedan körs från början.

\begin{asmcode}
    ldi r16, 0xFF
    out DDRB, r16
    ldi r16, 1
    ldi r17, 4
loop:
    out PORTB, r16
    lsl r16
    dec r17
    brne loop
\end{asmcode}

\begin{enumerate}
  \item Hur många gånger körs loopen, och vilket värde har \code{PORTB} när loopen är klar?
    \worth{1}
  \item Hur många klockcykler tar loopen, från den första \code{out} till och med den sista
    \code{brne}, och hur lång tid är det? \worth{1}
\end{enumerate}

\qpart{c} En knapp är kopplad mellan PD3 och jord, och en lysdiod med förkopplingsmotstånd mellan
PB0 och jord. Skriv ett komplett program, med initiering, där lysdioden lyser så länge knappen
hålls nedtryckt och är släckt annars. Använd stiftets inbyggda pull-up-motstånd. \worth{3}
